\begin{lemma}[elliptic values and specialization]\label{lem:elliptic-specialization}
Let $N\in\mathcal N$ and let
$h_N(T)=\prod_{\tau}(T-T_N(\tau))$, where $\tau$ runs over the
inequivalent elliptic points of $\Gamma_0(N)^+$.
The fifteen polynomials and their effective indices are given in
Table~\ref{tab:elliptic-values}.
For every prime $p\ge5$ with $p\ne N$, the reduction
$\overline h_N$ is squarefree, and its roots in $\Fb$ are exactly the
values with effective index $e>1$. Consequently,
\[
\gcd(\mathcal S_{N,p},\overline h_N)=\prod_{e_i>1}(T-t_i).
\]
\end{lemma}
\begin{proof}
The characteristic-zero polynomials are those of
\citet[\S4.2, Table~1]{JST19}, with $T_N=q^{-1}+O(q)$.
Their elliptic orders give the effective indices of
Lemma~\ref{lem:inertia}. For $N\in\mathcal N_0$, the Fricke fixed values
are given by $u_N^2=b_N$. At the remaining levels, the smooth models
$v^2=f_N(u)$ of Section~\ref{sec:comp} have involution
$(u,v)\mapsto(u,-v)$ and coordinate $T_N=u-s_N$;
the branch values are the roots of $f_N(T+s_N)$, all of index two.
The additional factors are $T+2$ of index two at $N=17,29$,
$T$ of index two at $N=41$, $T+3$ of index three at $N=19$,
and $T$ of index three at $N=31$. These are obtained from the
elliptic fibres over $j=1728$ and $j=0$.

We justify specialization of the support. An effective stabilizer arises
from an automorphism of $(E,C)$ or a Fricke fixed point. For $p\ge5$,
nontrivial reduced automorphisms occur only at $j=0,1728$; their
prime-to-$p$ automorphisms lift with the corresponding CM elliptic curve,
and an invariant prime-to-$p$ cyclic subgroup lifts uniquely over a
strict henselian base. Conversely, these automorphisms persist on reduction.
The displayed Fricke branch equations give finite \'etale fixed-point
loci away from $6N$: this follows directly from their discriminants.
Thus the listed elliptic values specialize to precisely the values with
$e>1$.

Exact discriminant calculation for every displayed $h_N$ gives
\[
\operatorname{Disc}(h_N)=\pm2^{\lambda_N}3^{\mu_N}N^{\nu_N}
\qquad(\lambda_N,\mu_N,\nu_N\ge0).
\]
The exact values and model calculations are supplied in
\path{Fricke_Elliptic_Values_Fifteen_Levels_v1_0}.
Hence distinct elliptic values cannot collide at any admissible prime,
so $\overline h_N$ is squarefree. Intersecting its root set with the
supersingular locus proves the gcd formula.
\end{proof}

\begin{table}[p]\centering\small
\setlength{\tabcolsep}{4pt}
\begin{tabular}{cll}\toprule
$N$ & $h_N(T)$ & Effective indices\\\midrule
2 & $(T+104)(T-152)$ & $4,\ 2$\\[3pt]
3 & $(T+42)(T-66)$ & $6,\ 2$\\[3pt]
5 & $(T+16)(T^{2}-12T-464)$ & $2,\ 2$\\[3pt]
7 & $(T+10)(T+9)(T-18)$ & $2,\ 3,\ 2$\\[3pt]
11 & $(T+6)(T^{3}-2T^{2}-76T-212)$ & $2,\ 2$\\[3pt]
13 & $(T+4)(T+3)(T^{2}-4T-48)$ & $2,\ 3,\ 2$\\[3pt]
17 & $(T+2)(T^{4}+2T^{3}-39T^{2}-176T-212)$ & $2,\ 2$\\[3pt]
19 & $(T+4)(T+3)(T^{3}-4T^{2}-16T-12)$ & $2,\ 3,\ 2$\\[3pt]
23 & $\begin{aligned}[t]&(T^{3}-2T^{2}-17T-25)\\[-1pt]&{}\cdot(T^{3}+6T^{2}+11T+7)\end{aligned}$ & $2,\ 2$\\[3pt]
29 & $(T+2)(T^{6}+2T^{5}-17T^{4}-66T^{3}-83T^{2}-32T-4)$ & $2,\ 2$\\[3pt]
31 & $T(T^{3}-17T-27)(T^{3}+4T^{2}+3T+1)$ & $3,\ 2,\ 2$\\[3pt]
41 & $\begin{aligned}[t]T\bigl(&T^8+4T^7-8T^6-66T^5-120T^4\\[-1pt]&{}-56T^3+53T^2+36T-16\bigr)\end{aligned}$ & $2,\ 2$\\[3pt]
47 & $\begin{aligned}[t]&(T^{5}-5T^{3}-20T^{2}-24T-19)\\[-1pt]&{}\cdot(T^{5}+4T^{4}+7T^{3}+8T^{2}+4T+1)\end{aligned}$ & $2,\ 2$\\[3pt]
59 & $\begin{aligned}[t]&(T^{3}+2T^{2}+1)\\[-1pt]&{}\cdot(T^{9}+2T^{8}-4T^{7}-21T^{6}-44T^{5}-60T^{4}-61T^{3}-46T^{2}-24T-11)\end{aligned}$ & $2,\ 2$\\[3pt]
71 & $\begin{aligned}[t]&(T^{7}-7T^{5}-11T^{4}+5T^{3}+18T^{2}+4T-11)\\[-1pt]&{}\cdot(T^{7}+4T^{6}+5T^{5}+T^{4}-3T^{3}-2T^{2}+1)\end{aligned}$ & $2,\ 2$\\[3pt]
\bottomrule\end{tabular}
\caption{Elliptic-value polynomials at all fifteen prime levels, in the
normalization $T_N=q^{-1}+O(q)$, from \citet[\S4.2, Table~1]{JST19}.
The indices follow the displayed factors; each index applies to every root
of its factor. At admissible primes, the supersingular roots of
$\overline h_N$ are exactly the thick values removed in
\eqref{eq:elliptic-gcd}.}
\label{tab:elliptic-values}
\end{table}
